% Lösungen Einheit 3 von 5 – Brüche vergleichen (zusammenbau v0.1)
\einheitenkopf{Einheit 3 von 5 · Brüche vergleichen}

%% TODO Lösung der Erklärzeile 18g) fehlt in der Bank
\erg{18}{a) gleicher Nenner \quad b) $\frac{7}{11}$ \quad c) $\frac{7}{12}$ \quad d) $\frac{5}{6}$ \quad e) $\frac{11}{15}$ \quad f) $\frac{13}{20}$ \quad g) \ldots \quad h) $\frac{3}{5}$ \quad i) $\frac{5}{8}$, denn $\frac{4}{9} < \frac{1}{2} < \frac{5}{8}$ \quad j) $\frac{3}{4}$, denn $\frac{21}{28} > \frac{20}{28}$ \quad k) $\frac{3}{8}$ am dritten, $\frac{7}{8}$ am siebten Teilstrich nach der $0$ \quad l) $\frac{1}{6}$, $\frac{2}{3}$, $\frac{5}{6}$ \quad m) z. B. $0{,}5$, denn $\frac{1}{4} = 0{,}25$ und $\frac{3}{5} = 0{,}6$; jede Zahl dazwischen ist richtig}

18k)

\zahlenstrahl[xmin=0,xmax=1,xstep=0.125,karo=1.2]{0/0, 1/1, 0.375/\frac{3}{8}, 0.875/\frac{7}{8}}

\erg{19}{a) z. B. $\frac{3}{10}$ (erweitert: $\frac{2}{10}$ und $\frac{4}{10}$)}
\erg{20}{a) $\frac{5}{6}$ ist größer: Bei $\frac{2}{3}$ fehlt ein Drittel, bei $\frac{5}{6}$ nur ein Sechstel – die fehlenden Stücke sind verschieden groß. \quad b) Beim Fünftel wird das Ganze in weniger Teile geteilt, jedes Teil ist also größer. \quad c) $\frac{3}{4}$ \quad d) Emil: $\frac{2}{3} = \frac{16}{24} > \frac{15}{24} = \frac{5}{8}$}
